\documentclass{article}

\usepackage{amsmath,amssymb}
\usepackage{xurl}
\usepackage[hidelinks]{hyperref}
\setlength{\emergencystretch}{2em}

\input{command}

\title{Wolf Chapter~5 Special-Function Jensen and Lieb Concavity}
\author{}
\date{2026-08-24}

\begin{document}
\maketitle
\gapnote{false-source}{resolved}

\section{Scope}

This note records two resolved parts of Chapter~5 of
Ref.~\cite{Wolf2012Quantum}. First, the formal development proves the
special-function Jensen inequalities corresponding to Wolf's Corollary~5.2,
after correcting its false logarithmic clause. Second, it proves the
finite-dimensional Ando--Lieb theorem throughout the exponent range stated in
Wolf's Theorem~5.15. The separate scope gap for Wolf's arbitrary
operator-convex and operator-monotone functions is recorded in
\path{docs/paper-gaps/wolf_ch5_operator_convexity_scope.tex}.

All operators act on a finite-dimensional Hilbert space, identified with
$D\times D$ complex matrices. The relation $A\leq B$ is the Loewner order. A
linear map $T$ is positive if it preserves positive semidefiniteness,
subunital if $T(\Id)\leq\Id$, and unital if $T(\Id)=\Id$. For $A\geq 0$, the
power $A^p$ is defined by continuous functional calculus; for $A>0$, the same
applies to $\log A$.

The relevant formal declarations are proved theorems rather than
axioms.\footnote{Formal declarations:
\leanid{posMap_rpow_concave_jensen},
\leanid{posMap_rpow_convex_jensen},
\leanid{posMap_log_concave_jensen}, and
\leanid{lieb_concavity_posDef} in
\doclink{QICLean/Analysis/LiebConcavity}.}

\section{Positive-map Jensen inequalities}

\subsection{The assertions}
\label{ssec:wolf_ch5_operator_jensen_lieb_assertions}

For a positive subunital map $T$ and a positive semidefinite matrix $A$, the
development proves
\begin{align}
    T(A^p)&\leq(T(A))^p,
    &
    0&\leq p\leq 1,
    \label{eq:wolf_ch5_operator_jensen_lieb_rpow_concave_jensen}\\
    (T(A))^p&\leq T(A^p),
    &
    1&\leq p\leq 2.
    \label{eq:wolf_ch5_operator_jensen_lieb_rpow_convex_jensen}
\end{align}
For a positive unital map $T$ and $A>0$, it proves
\begin{align}
    T(\log A)&\leq\log(T(A)).
    \label{eq:wolf_ch5_operator_jensen_lieb_log_concave_jensen}
\end{align}
These are the operator-valued Jensen inequalities associated with Wolf's
Corollary~5.2 \cite[Corollary~5.2]{Wolf2012Quantum}.\footnote{Map-facing
interfaces:
\leanid{IsPositiveMap.rpow_concave_jensen},
\leanid{IsPositiveMap.rpow_convex_jensen}, and
\leanid{IsPositiveMap.log_concave_jensen} in
\doclink{QICLean/Channel/Schwarz/OperatorConvexity}.}

Equation~\ref{eq:wolf_ch5_operator_jensen_lieb_rpow_concave_jensen} expresses
operator concavity of $x\mapsto x^p$ for $p\in[0,1]$.
Equation~\ref{eq:wolf_ch5_operator_jensen_lieb_rpow_convex_jensen} expresses
operator convexity for $p\in[1,2]$.
Equation~\ref{eq:wolf_ch5_operator_jensen_lieb_log_concave_jensen} expresses
operator concavity of the logarithm. The logarithmic theorem requires
unitality; Wolf's common subunital hypothesis is insufficient.

\subsection{The special-function proof route}

The formal library supplies integral representations for real powers, including
\leanid{CFC.exists_measure_nnrpow_eq_integral_cfcₙ_rpowIntegrand₀₁} and its
$[1,2]$ companion in
\path{Mathlib/Analysis/SpecialFunctions/ContinuousFunctionalCalculus/Rpow/IntegralRepresentation.lean}.
It also supplies operator concavity of $x\mapsto x^p$ on $[0,1]$ through
\leanid{CFC.concaveOn_rpow}, operator concavity of the logarithm through
\leanid{CFC.concaveOn_log}, and the required Bochner-integral convexity
transfer for C\textsuperscript{*}-algebras. Operator convexity of
$x\mapsto x^p$ on $[1,2]$ is recorded as \leanid{CFC.convexOn_rpow} in
\doclink{QICLean/Analysis/RpowConvexity}.

A general passage from operator concavity of $f$ to
\begin{align}
    T(f(A))&\leq f(T(A))
    \label{eq:wolf_ch5_operator_jensen_lieb_general_concave_jensen_form}
\end{align}
for positive maps is the positive-map operator Jensen principle used in
Wolf's Theorems~5.10--5.13 \cite[Theorems~5.10--5.13]{Wolf2012Quantum}.
Mathlib~4.31 contains no general theorem of this scope for positive maps.
Instead, the concave real-power case is proved by the L\"owner-integrand
theorem
\leanid{TNLean.OperatorJensen.positiveMap_rpowIntegrand₀₁_jensen} and ordered
Bochner integration.

The logarithmic theorem follows by taking the right limit $p\to 0^+$ in the
concave real-power inequality and using
\leanid{CFC.tendsto_cfc_rpow_sub_one_log}. The convex real-power case uses the
same positive-map Jensen method with the convex integrand for $p\in[1,2]$.

\subsection{Bare positivity is sufficient}
\label{ssec:wolf_ch5_operator_jensen_lieb_bare_positivity}

The source statements assume that $T$ is positive, not completely positive
\cite[Corollary~5.2 and Theorems~5.10--5.13]{Wolf2012Quantum}. A proof that
required complete positivity would therefore strengthen the source
hypotheses.

Wolf avoids this strengthening by restricting $T$ to the abelian
C\textsuperscript{*}-algebra generated by $A$ and $\Id$. On this algebra every
positive map is automatically completely positive, and the restricted map has
a completely positive extension to the full matrix algebra
\cite[Propositions~1.6--1.7 and Theorems~5.10--5.13]{Wolf2012Quantum}. Thus
the source proof reduces the bare positive-map case to a completely positive
map without changing the theorem's hypotheses.

The formal finite-POVM scaffold implements the corresponding matrix
reduction.\footnote{Formal development:
\doclink{QICLean/Channel/Schwarz/OperatorJensenAux}, containing finite-POVM
compression, the Schur-complement inverse bound, and the resolvent inequality.}
Let $P_1,\ldots,P_n$ be the spectral projections of $A$. Positivity of $T$
gives positive matrices $T(P_i)$, and linearity gives
\begin{align}
    \sum_{i=1}^nT(P_i)
        &=T\!\left(\sum_{i=1}^nP_i\right)
    \label{eq:wolf_ch5_operator_jensen_lieb_povm_sum_first}\\
        &=T(\Id).
    \label{eq:wolf_ch5_operator_jensen_lieb_povm_sum_second}
\end{align}
This is the identity used in Wolf's block-positivity argument
\cite[Theorem~5.12]{Wolf2012Quantum}. The construction uses a positive family
and a defect block, not a global Kraus or Stinespring representation of $T$.
It therefore applies under \leanid{IsPositiveMap} and subunitality.

The pointwise finite-POVM resolvent inequality
\leanid{povm_resolvent_inv_le} is integrated through the real-power
representation to prove
\leanid{posMap_rpow_concave_jensen}. Almost-everywhere Loewner inequalities
are integrated using \leanid{integral_mono_ae}. The positive-cone
specialization is
\leanid{TNLean.OperatorJensen.integral_nonneg_matrix_of_ae}, which states that
the Bochner integral of an almost-everywhere positive semidefinite
matrix-valued function is positive semidefinite.

The logarithmic theorem
\leanid{posMap_log_concave_jensen} then follows by applying
\leanid{CFC.tendsto_cfc_rpow_sub_one_log} to both $A$ and $T(A)$ and using
closedness of the finite-dimensional Loewner cone. The convex theorem
\leanid{posMap_rpow_convex_jensen} uses the corresponding convex integrand and
the input \leanid{CFC.convexOn_rpow}.

\subsection{The false logarithmic clause}

Wolf's Corollary~5.2 prints the logarithmic inequality under the same
subunital hypothesis as the real-power inequalities:
\begin{align}
    T(\log A)&\leq\log(T(A)),
    &
    A&>0.
    \label{eq:wolf_ch5_operator_jensen_lieb_printed_subunital_logarithm}
\end{align}
This is item~3 of Corollary~5.2 in Ref.~\cite{Wolf2012Quantum}. It is false.

In dimension one, choose $0<c<1$, define $T_c(x)=cx$, and set $A=1$. Then
\begin{align}
    T_c(1)&=c\leq 1,
    \label{eq:wolf_ch5_operator_jensen_lieb_scalar_map_subunital}\\
    T_c(\log 1)&=0,
    \label{eq:wolf_ch5_operator_jensen_lieb_scalar_logarithm_left}\\
    \log(T_c(1))&=\log c<0.
    \label{eq:wolf_ch5_operator_jensen_lieb_scalar_logarithm_right}
\end{align}
Consequently,
Eq.~\ref{eq:wolf_ch5_operator_jensen_lieb_printed_subunital_logarithm} would
assert $0\leq\log c$, a contradiction. The corrected declaration
\leanid{IsPositiveMap.log_concave_jensen} assumes
\begin{align}
    T(\Id)&=\Id
    \label{eq:wolf_ch5_operator_jensen_lieb_corrected_logarithm_unitality}
\end{align}
and proves
Eq.~\ref{eq:wolf_ch5_operator_jensen_lieb_log_concave_jensen}. This is a
source correction rather than a formalization-only restriction.

The real-power theorems retain Wolf's bare subunital hypothesis. Former
pass-through declarations that merely repackaged them in the exact exponent
parameterization of items~1--2 were removed as unused.

\subsection{Jensen verdict}

Both real-power Jensen inequalities are proved under the source hypotheses.
The logarithmic inequality is proved under the necessary unital correction.
The concave power proof uses finite-POVM compression, the pointwise resolvent
inequality, and ordered Bochner integration. The logarithmic proof takes the
right limit of the concave power theorem. The convex proof uses the
corresponding convex integrand. The required integral representations are
present.

This verdict concerns only these special functions. It does not resolve the
separate scope question for arbitrary operator-convex and operator-monotone
functions.

\section{The Lieb concavity theorem}
\label{sec:wolf_ch5_operator_jensen_lieb_lieb}

\subsection{The assertion}

Let $s\in[0,1]$, let $K$ be any complex matrix, and let $A,B>0$. Define
\begin{align}
    F_s(A,B)
        &=\tr\!\left(K^\dagger A^sKB^{1-s}\right).
    \label{eq:wolf_ch5_operator_jensen_lieb_boundary_functional}
\end{align}
The formal theorem \leanid{lieb_concavity_posDef}, whose input hypotheses use
\leanid{Matrix.PosDef}, states that $F_s$ is jointly concave on
positive-definite pairs. This is the boundary-exponent case of the Ando--Lieb
functional in Wolf's Theorem~5.15
\cite[Theorem~5.15]{Wolf2012Quantum}.

The declaration \leanid{lieb_concavity_psd}, whose input hypotheses use
\leanid{Matrix.PosSemidef}, extends this theorem to positive-semidefinite
inputs. The declaration \leanid{lieb_concavity_subboundary} then proves the
full exponent range
\begin{align}
    x&\geq 0,
    &
    y&\geq 0,
    &
    x+y&\leq 1
    \label{eq:wolf_ch5_operator_jensen_lieb_full_exponent_region}
\end{align}
stated in Wolf's Theorem~5.15
\cite[Theorem~5.15]{Wolf2012Quantum}.

For $r=x+y>0$, set
\begin{align}
    r&=x+y,
    &
    s&=\frac{x}{r}.
    \label{eq:wolf_ch5_operator_jensen_lieb_subboundary_parameters}
\end{align}
Then
\begin{align}
    A^x&=(A^r)^s,
    \label{eq:wolf_ch5_operator_jensen_lieb_subboundary_power_a}\\
    B^y&=(B^r)^{1-s}.
    \label{eq:wolf_ch5_operator_jensen_lieb_subboundary_power_b}
\end{align}
The proof combines operator concavity of $C\mapsto C^r$, monotonicity of the
boundary trace functional, and \leanid{lieb_concavity_psd}. When $r=0$, the
functional is constant.

For positive-semidefinite inputs, the proof replaces each matrix by
\begin{align}
    A_{i,\varepsilon}&=A_i+\varepsilon\Id,
    &
    B_{i,\varepsilon}&=B_i+\varepsilon\Id,
    &
    \varepsilon&>0.
    \label{eq:wolf_ch5_operator_jensen_lieb_regularized_inputs}
\end{align}
These matrices are positive definite. Applying the positive-definite theorem
and taking $\varepsilon\to 0^+$ proves \leanid{lieb_concavity_psd}. The powers
converge because $t\mapsto t^s$ is continuous on the nonnegative real axis for
$0<s<1$, including at zero, and the trace functional is continuous. No
logarithmic singularity occurs in this regularization.

\subsection{The integral representation}
\label{ssec:wolf_ch5_operator_jensen_lieb_integral_route}

A direct matrix resolvent formula for $A^sB^{1-s}$ is valid only when $A$ and
$B$ commute. For noncommuting matrices, $A^sB^{1-s}$ need not be
self-adjoint, so that formula cannot prove the general Lieb theorem.

The correct representation uses left and right multiplication operators on
$\MD$. Define
\begin{align}
    L_A(X)&=AX,
    &
    R_B(X)&=XB.
    \label{eq:wolf_ch5_operator_jensen_lieb_left_right_multiplication}
\end{align}
These operators commute:
\begin{align}
    L_AR_B(X)
        &=AXB
    \label{eq:wolf_ch5_operator_jensen_lieb_superoperator_commutation_first}\\
        &=R_BL_A(X).
    \label{eq:wolf_ch5_operator_jensen_lieb_superoperator_commutation_second}
\end{align}
For $A,B\geq 0$, they are positive operators on the Hilbert--Schmidt space
with inner product
\begin{align}
    \langle X,Y\rangle_{\mathrm{HS}}
        &=\tr(X^\dagger Y).
    \label{eq:wolf_ch5_operator_jensen_lieb_hilbert_schmidt_inner_product}
\end{align}
The Lieb functional becomes
\begin{align}
    \tr\!\left(K^\dagger A^sKB^{1-s}\right)
        &=
        \left\langle
            K,L_A^sR_B^{1-s}K
        \right\rangle_{\mathrm{HS}}.
    \label{eq:wolf_ch5_operator_jensen_lieb_superoperator_pairing}
\end{align}
For $0<s<1$, the commuting-superoperator integral representation is
\begin{align}
    L_A^sR_B^{1-s}
        &=
        \frac{\sin(\pi s)}{\pi}
        \int_0^\infty
            t^{s-1}
            L_A(L_A+tR_B)^{-1}R_B
        \,dt.
    \label{eq:wolf_ch5_operator_jensen_lieb_superoperator_integral}
\end{align}
Joint concavity of the integrand in $(A,B)$ for each fixed $t$ follows through
the corresponding resolvent argument. Pairing with $K$ in the
Hilbert--Schmidt inner product, integrating, and using linearity then proves
joint concavity.

Mathlib~4.31 contains neither this superoperator integral representation nor
the required matrix-mean and resolvent-monotonicity interface for the commuting
pair $(L_A,R_B)$. These ingredients were developed within the formalization.
The proof proceeds through the scalar boundary integral, the operator integral
representation on the Kronecker model, concavity of the resolvent integrand,
and the resulting operator concavity.

\subsection{Lieb verdict}

The Lieb concavity assertion in Wolf's Theorem~5.15 is correct
\cite[Theorem~5.15]{Wolf2012Quantum}. The formal proof uses the commuting
superoperators $L_A$ and $R_B$, not a matrix resolvent identity that would
require $A$ and $B$ to commute.

The declaration \leanid{lieb_concavity_posDef} proves the positive-definite
boundary-exponent theorem. The declaration \leanid{lieb_concavity_psd} removes
the positive-definiteness restriction by regularization, and
\leanid{lieb_concavity_subboundary} proves the full region in
Eq.~\ref{eq:wolf_ch5_operator_jensen_lieb_full_exponent_region}. Three
downstream corollaries, including the $K=\Id$ case and separate concavity in
each argument, build on the positive-definite theorem and are now axiom-free.
The former positive-semidefinite and sub-boundary restrictions have both been
eliminated.

\section{Conclusion}

The special-function Jensen theorems and the finite-dimensional Ando--Lieb
theorem are proved. The real-power Jensen inequalities retain Wolf's positive
subunital hypotheses. The logarithmic clause in Corollary~5.2 of
Ref.~\cite{Wolf2012Quantum} is false under subunitality and is corrected by
requiring unitality. The Lieb theorem now covers positive-semidefinite inputs
and all exponents $x,y\geq 0$ with $x+y\leq 1$.

The relevant functional-calculus inputs are
\leanid{CFC.exists_measure_nnrpow_eq_integral_cfcₙ_rpowIntegrand₀₁},
\leanid{CFC.concaveOn_rpow}, \leanid{CFC.convexOn_rpow}, and
\leanid{CFC.concaveOn_log}. The three positive-map Jensen inequalities and the
Lieb theorem treated in this note are proved.

\bibliographystyle{alpha}
\bibliography{references}

\end{document}
